\input{../tex-inputs/latex-leseansicht-vorspann}

\section[Arthur Schnitzler an Gustav Schwarzkopf, 5. 10. 1904]{L04028 Arthur Schnitzler an Gustav Schwarzkopf, 5. 10. 1904}
\nopagebreak\mylabel{L04028v}
\rehead{ }\normalsize\beginnumbering\briefempfaengerindex{Schwarzkopf, Gustav@\textsc{Schwarzkopf, Gustav}!zzzSchnitzler, Arthur@\emph{von Arthur Schnitzler}!1904-10-051@{5. 10. 1904}|(be}
\toendnotes[C]{\smallbreak\pagebreak[2]}
\correspDesc{Versand  durch Arthur Schnitzler am 5. 10. 1904 in Wien
\newline{}Erhalt  durch Gustav Schwarzkopf am 5. 10. 1904 in Wien}\toendnotes[C]{\smallbreak}
\Standort{CUL, Schnitzler, B 96.}
\physDesc{Brief, 1 Blatt, 1 Seite, 150 Zeichen
\newline{}Handschrift: Bleistift, deutsche Kurrent}\toendnotes[C]{\smallbreak}
\pstart
           \raggedleft{}{\pb}5. X. 904\pend
           
\pstart{}lieber Guſtav\pend\vspace{0.5em}
\pstart
           bitte nehmen Sie freundlich dieſen Sitz\eventindex{Theater in der Josefstadt@\textbf{Theater in der Josefstadt}!Aufführung von Herzogin Crevette, 5.10.1904@Aufführung von Herzogin Crevette, 5.10.1904|pwv} an – Olga\pwindex{Schnitzler, Olga 17.\,1.\,1882 Wien – 13.\,1.\,1970 Lugano@\textsc{Schnitzler, Olga} (17.\,1.\,1882 Wien – 13.\,1.\,1970 Lugano), \emph{Schauspielerin, Sängerin}|pw} iſt nicht ganz wohl u ka{\geminationn} daher nicht
               gehn. Auf Wiederſehn heut Abend\pend
           
\pstart
           Herzlich{\\[\baselineskip]} Ihr{\\[\baselineskip]}\spacefill\mbox{A.}\pend
           \leftskip=0em{}\selectlanguage{ngerman}\endnumbering\briefempfaengerindex{Schwarzkopf, Gustav@\textsc{Schwarzkopf, Gustav}!zzzSchnitzler, Arthur@\emph{von Arthur Schnitzler}!1904-10-051@{5. 10. 1904}|)be}\mylabel{L04028h}
\begin{anhang}
\end{anhang}\newcommand{\dateiname}{L04028}\newcommand{\titel}{Arthur Schnitzler an Gustav Schwarzkopf, 5. 10. 1904}\newcommand{\editorInnen}{Herausgegeben von Jahnke, SelmaMüller, Martin Anton}\input{../tex-inputs/latex-leseansicht-abspann}
